\section{Experimental Development and Scope Boundaries}
\label{sec:appendix_development}

The reported system was developed through several experiments with different input
contracts. The following account identifies the decisions that connect the main results;
the released records retain the detailed configurations, failures, and intermediate
outputs. Historical scores are not merged into the current v5 leaderboard.

\textbf{Early replication and paired-input prototypes.} Early work included direct
answering, MedRGAG reproduction and local adaptations, a world-model proxy, and
Delta/Residual Adapter variants. The apparent advantage of world-model sorting did not
survive a matched prompt-and-input comparison: both it and blind direct sorting scored
64/85. A subsequent counterfactual experiment identified a label-oracle condition;
that condition was rejected as valid method evidence, and the corrected label-free
version did not establish the proposed mechanism. A later Residual Adapter achieved
80.0\% with MA-RAG on a 500-pair four-choice task, while a separately trained
without-MA configuration reached 81.8\%. Those experiments provided useful controls,
but their paired-input and constructed-choice permissions differ from the current
independent-endpoint native protocol. They neither establish synergy nor supply extra
v5 baseline scores.

\textbf{Baseline lineage.} The evaluation progressed from earlier source-level and
v2/v3 comparisons to the frozen v5 selection. M1 denotes MedRGAG's initial retrieval
control, not a separate MedRAG family; its earlier completed evaluations do not provide
a full v5 result. Historical M6 denoted i-MedRAG. Its original v5 queue paused after
52 of 13,905 inputs; a later 128-input performance check was not a completed v5
evaluation. The selected current M6 is MedGENIE and preserves a different method
identity. M4's structured-output repair and M7's typed-control revision are included
in the selected baseline versions, with earlier failures retained in their own records.

\textbf{From errors to terminal heads.} R2 development first tested ordinary revision,
free-form support audits, local withdrawal, atomic decisions, additional guideline
context, and longer reasoning. Their failures led to a narrower division of labor:
the model extracts current facts, and a fixed rule executor handles conditions,
exceptions, question scope, and independent support paths. An answer-assembly type
error was repaired using saved outputs, and its earlier negative conclusion was
withdrawn rather than retained as a method failure. The final default preserves
UNKNOWN initial choices and records its remaining-support regressions.

R3 first tested evidence audits, query changes, and training-case memory. A high-scoring
nearest-neighbor head was withdrawn from the core method because it returned retrieved
training labels instead of making the intended original-model decision. The subsequent
public-catalog head is a distinct implementation. Patient presentation, catalog
selection, permutations, and prior correction were evaluated separately. More elaborate
pairwise comparisons, observation selection, knowledge allocation, relation counting,
and NLI candidates did not replace the retained head. R1 then extended the frozen
diagnostic and rule components beyond R3, with complete seven-system accounting.
Single-choice extension and revised fact extraction had negative results; a relation
candidate remained a separate two-configuration experiment.

\textbf{R4: conditional consequence execution.} The original 13-unit HIV stratum has
seven-system baseline results, but subsequent R4 head development used two backbones.
Clinical prompt and evidence-binding variants did not establish a retained clinical
head. Formal and biological tasks informed the later implementation, whose adopted
scope is explicitly supplied physiological models and states. Frozen validation on
24 new virtual profiles, each with 12 questions, is reported in
Appendix~\ref{sec:appendix_model}. It is completed conditional model validation, not
seven-system optimization of the 13 HIV units or evidence of unrestricted clinical
causal reasoning.

\textbf{R5: retained readout and unadopted learning candidates.} The historical F
component combines two NLI paths and two label mappings, or up to five native-answer
draws for its other supported tasks. Its 208-unit development/validation panel contains
64 initial development units, 128 extension units, and 16 new groups evaluated after
freezing. Across the 208 units, weighted paired correctness changes from 60.19\% to
65.90\% for Llama and from 59.17\% to 66.47\% for Qwen. Within the 16 frozen-after
groups, both-correct pairs improve from 11 to 12 for Llama and from 12 to 13 for Qwen
relative to matched single answers. These
limited results support retaining F as a component. Later support audits, learned
representations, and fallback gates failed their controls and were not adopted; a
further five-head training proposal remained unexecuted. They are not extra successful
R5 methods hidden within the final composition.

\textbf{Composition and verification follow-up.} CF-MoA subsequently tested global
aggregation and additional support, comparison, and maintenance variants. These were
executed experiments with negative or mixed results, not merely unimplemented plans.
A minimal wrapper then preserved the established B operations, while ordinary fresh
answer generation and finally candidate verification were tested as additional
interventions. Earlier 121-input and 52-input panels informed that process. The
seven-system component migration replay demonstrated equivalence of saved results;
its cached callbacks were not new quality measurements. The final six-arm, full-v5
experiment applies only to M4 and M5. Its 162,979-request physical ledger covers that
run, not the total cost of the development history summarized here.
